\subsection{Scenario generation}

The scenario generator creates the demand used by the rest of the pipeline. In
this work, the demand must satisfy four constraints:
\begin{enumerate}
\item \textbf{Network-aware}: Origins and destinations are sampled from actual
      SUMO road edges.
\item \textbf{Zone-aware}: Demand is expressed at the TAZ level.
\item \textbf{Temporally structured}: The daily profile includes morning and
      evening peaks.
\item \textbf{Directional}: Peak periods can concentrate trips from
      residential zones toward employment zones and back.
\end{enumerate}

This module is thesis code. It does not claim to estimate real Brussels demand
from counts, surveys or mobile-phone traces. Instead, it constructs a
controlled synthetic demand that is consistent with the available network
files and is rich enough to test the detour-generation and evaluation
pipeline. The standard modelling components used for this purpose are
identified below, and the thesis-specific choices are the proxies,
normalisations and sampling procedures used to turn them into reproducible
SUMO inputs.

\subsubsection{Module overview}
The generation is split into three stages:
\begin{enumerate}
  \item \textbf{Stage 1 - Demand generation}: construct 24 hourly
        Origin\nobreakdash–Destination (OD) matrices using a gravity model \cite{gravity_model}
        with an optional directional commute overlay.
  \item \textbf{Stage 2 - Route library}: build a small set of candidate
        routes for each active OD pair using Yen's $k$-shortest
        paths algorithm \cite{yen_algorithm}.
  \item \textbf{Stage 3 - Trip instantiation}: assign a departure time
        and valid endpoint edges to every trip. The full route is assigned
        later, inside the Monte Carlo stage.
\end{enumerate}

\subsubsection{Stage 1 : Demand generation}
\label{ssec:stage_1_demand_generation}
Stage 1 produces hourly OD matrices, not individual routes. The output of this
stage is a set of zone-to-zone trip counts for each hour of the day. Individual
departure times and SUMO edge endpoints are assigned only in Stage 3.

\paragraph{Traffic Analysis Zones}
The spatial unit of the model is the \emph{Traffic Analysis Zone} (TAZ),
a standard construct in transport demand modelling that aggregates individual
trip origins and destinations into discrete geographic areas.  
\href{https://sumo.dlr.de/docs/Contributed/SUMOPy/Demand/Zone_To_Zone.html}{TAZs in SUMO} 
are defined as polygons that can be associated with the road network.
Each TAZ is defined by the set of SUMO road edges it contains.

A scalar \emph{zone mass} $m_i$ is assigned to each TAZ as the capacity of its drivable
edge lengths : 
\begin{equation}
  m_i = \sum_{e \in \mathcal{E}_i} \ell(e) \cdot n(e)
  \label{eq:zone_mass}
\end{equation}
where $\ell(e)$ is the length of edge $e$ in meters and $n(e)$ is its lane count and 
$\mathcal{E}_i$ is the set of edges in TAZ $i$. 

Equation~\eqref{eq:zone_mass} should therefore be interpreted as an infrastructure-based 
proxy for the relative importance of each TAZ, rather than as a calibrated measure of 
real travel demand. The formula only accounts for the total drivable road space inside a 
zone, through edge lengths and lane counts. The underlying assumption is that zones containing 
more and larger road infrastructure are likely to support more traffic activity and can 
therefore be assigned a larger synthetic production or attraction mass. This approximation 
is used because no calibrated socio-economic productions or attractions, such as population, 
employment, land use, or observed trip data are available for this synthetic setting. In a 
classical four-step transport model, these quantities would instead be estimated from socio-economic 
data \cite{modelling_transport}. Here, lane
kilometres\footnote{A lane kilometre is one kilometre of road with one lane.}
are used as a practical surrogate: they are directly available from the
network and major road corridors generally carry more traffic than minor
links. If observed productions and attractions become available, they should
replace this proxy.

\paragraph{Temporal Demand Profile}
\label{par:temporal_demand_profile}

The total daily demand $N$ is distributed across the 24 hours of the
simulation using a \emph{two\nobreakdash-Gaussian profile}:
\begin{equation}
  p_h = \frac{a_{AM}\,g(h;\,\mu_{AM},\sigma) + a_{PM}\,g(h;\,\mu_{PM},\sigma)}
             {\displaystyle\sum_{h'=0}^{23}
              \bigl[a_{AM}\,g(h';\,\mu_{AM},\sigma)
                  + a_{PM}\,g(h';\,\mu_{PM},\sigma)\bigr]},
  \label{eq:twopeak_gaussian}
\end{equation}
where 
\begin{equation}
  g(h;\,\mu,\sigma) = \exp\!\left(-\frac{(h-\mu)^2}{2\sigma^2}\right),
\end{equation}
$\mu_{AM}$ and $\mu_{PM}$ are the morning and evening peak hours respectively and
$\sigma$ controls the temporal spread. The amplitudes $a_{AM}$ and $a_{PM}$ allow
the two peaks to differ in magnitude. The default configuration uses
$a_{AM} < a_{PM}$, representing a slightly larger afternoon peak, which is common in
urban networks \cite{gaussian_book}.

The continuous profile $p_h$ is converted into hourly trip counts $n_h$ with
the largest-remainder method. This keeps the total exactly equal to $N$ while
staying as close as possible to the continuous proportions.

Figure~\ref{fig:hourly_demand} illustrates the mathematical structure of the
two-peak profile before integer rounding.

\begin{figure}[H]
  \centering
  \safeincludegraphics[width=\textwidth]{illustrative_two_gaussian_profile.pdf}
  \caption{Illustrative two-Gaussian daily demand profile used to distribute
           the total daily demand across hours. The morning and evening
           components are centred on $\mu_{AM}$ and $\mu_{PM}$, and the final hourly
           share is proportional to their weighted sum
           $a_{AM} g(h;\mu_{AM},\sigma)+a_{PM} g(h;\mu_{PM},\sigma)$.}
  \label{fig:hourly_demand}
\end{figure}

\paragraph{Gravity OD model}
\label{par:gravity_od_model}

The gravity model itself is not a new model introduced by this thesis. It is a
standard transport-planning model used to distribute trips between origins and
destinations as a function of origin mass, destination mass and travel
impedance \cite{gravity_model,modelling_transport}. What is specific to this
work is its synthetic implementation: zone mass is approximated from lane
kilometres, impedance is estimated from sampled shortest paths between TAZs
and the resulting shares are converted into exact hourly integer OD matrices.

The spatial distribution of trips across TAZ pairs is governed by a
\emph{singly-constrained gravity model}. The \textbf{unnormalized} weight of
the OD pair $(i,j)$ is:
\begin{equation}
  w_{ij} = m_i \cdot m_j \cdot \exp(-\beta\, c_{ij}), \quad i \neq j,
  \label{eq:gravity_model}
\end{equation}
where $c_{ij}$ is the travel-cost impedance between zones $i$ and $j$ (shortest path length
in meters) and $\beta \geq 0$ is the distance decay parameter\footnote{A larger $\beta$
means trips decay faster with distance/cost (people are more reluctant to travel far), a smaller $\beta$
means the distribution is more spatially spread out. Wilson discusses this in the context of 
calibrating the model to observed data.\cite{gravity_model}}.

Pairs with infinite or negative impedance receive zero weight. The normalized
\emph{OD share matrix} is then:
\begin{equation}
  S_{ij} \;=\; \underbrace{\frac{m_i}{\sum_{k} m_k}}_{\text{origin share }O_i}
              \;\cdot\;
              \underbrace{\frac{m_j\,\exp(-\beta c_{ij})}
                               {\sum_{\ell} m_\ell\,\exp(-\beta c_{i\ell})}}_{\text{conditional destination probability }p(j\mid i)},
  \label{eq:od_share}
\end{equation}
so that, by construction, $\sum_j S_{ij} = m_i / \sum_k m_k$ and
$\sum_{ij} S_{ij} = 1$.

For each hour $h$, the share matrix is scaled to the corresponding integer
total $n_h$ with the largest-remainder method, giving an exact integer OD
matrix $T^h$ with
$\sum_{ij} T^h_{ij} = n_h$.

The model is \emph{singly-constrained on the production side}. In trip-distribution
terminology, productions are the trips generated by each origin zone, whereas
attractions are the trips received by each destination zone. A production-side
constraint therefore fixes the total fraction of trips that must leave each
origin. In Equation~\eqref{eq:od_share}, this production total is set by the
origin mass share $m_i / \sum_k m_k$: after the destination probabilities for
origin $i$ are normalised, they must sum to one, so the whole row must sum to
\begin{equation*}
  \sum_j S_{ij} \;=\; \frac{m_i}{\sum_k m_k}
\end{equation*}
regardless of the destinations available from that origin. Equivalently, in an
hour with $n_h$ trips, origin $i$ is assigned approximately
$n_h m_i / \sum_k m_k$ outgoing trips before integer rounding.

The destination side is different. The term $m_j\exp(-\beta c_{ij})$ still
affects how the trips from origin $i$ are distributed across destinations:
larger and closer destination zones receive higher conditional probability.
However, no target total is imposed on the number of trips ending in
destination $j$. The column sums $\sum_i S_{ij}$ are therefore left
unconstrained and emerge from the aggregation of all origin rows. This is why
the unnormalised weight $w_{ij} = m_i m_j \exp(-\beta c_{ij})$ can contain both
origin and destination masses while the final model is still only constrained
on the origin, or production, side.

In Wilson's balancing-factor formulation \cite{modelling_transport}, this
corresponds to enforcing only the origin-side factor $A_i$ and fixing the
destination-side factor $B_j \equiv 1$. A doubly-constrained model would
additionally enforce observed zonal attractions and would require both
balancing factors to be recovered by iterative proportional fitting. Since
zonal attractions are not observed in the present synthetic setting, enforcing
them would add parameters that cannot be identified from the available data.
For this reason, the singly-constrained formulation is used here: it relies
only on the information present in the network (zone masses and impedances)
and lets destination totals emerge from the destination masses and
distance-decay term $\exp(-\beta c_{ij})$.
Figure~\ref{fig:od_heatmap} gives a schematic view of the same directed OD
relations in matrix form.

\begin{figure}[H]
  \centering
  \safeincludegraphics[width=\textwidth]{illustrative_od_matrix.pdf}
  \caption{Schematic relationship between zone-to-zone trip relations and an
           OD matrix. Rows represent origin zones, columns represent
           destination zones, and each off-diagonal cell corresponds to the
           relative demand assigned to one directed OD pair. The diagonal is
           set to zero when intra-zonal trips are excluded.}
  \label{fig:od_heatmap}
\end{figure}

\paragraph{Commute Overlay (Optional)}
\label{par:commute_overlay}

When residential weights $r_i$ and employment weights $e_i$ are available,
the hourly OD shares can be tilted toward commuting patterns. Two additional
gravity matrices are built: one for the morning peak (AM), from residential
zones to employment zones and one for the evening peak (PM), in the reverse
direction:
\begin{align}
  w^{\mathrm{AM}}_{ij} &= m_i\, m_j\, f(c_{ij})\, r_i\, e_j, \\
  w^{\mathrm{PM}}_{ij} &= m_i\, m_j\, f(c_{ij})\, e_i\, r_j,
\end{align}
where $f(c_{ij}) = \exp(-\beta_c\, c_{ij})$ discounts longer or costlier
trips. In the AM matrix, $r_i$ and $e_j$ increase flows leaving residential
zones and arriving at employment zones. In the PM matrix, the same logic is
reversed.

The two directional matrices are blended with the base OD share matrix through
a time-varying mixture. For each hour $h$, the blended share is:
\begin{equation}
  S^h_{ij} \;\propto\; S^{\mathrm{base}}_{ij}
  + \lambda\, g(h;\,\mu_m,\,\sigma)\, S^{\mathrm{AM}}_{ij}
  + \lambda\, g(h;\,\mu_e,\,\sigma)\, S^{\mathrm{PM}}_{ij},
  \label{eq:commute_blend}
\end{equation}
where $g(h;\,\mu,\,\sigma)$ is a Gaussian kernel centred on hour $\mu$ with
spread $\sigma$. The parameter $\lambda$ (\texttt{peak\_weight}) controls how
strongly the commute pattern modifies the background gravity model. Setting
$\lambda = 0$ recovers the base model. After blending, the shares are
normalised before being converted into integer trips.
Figure~\ref{fig:commute_overlay} illustrates how the AM and PM terms tilt the
base gravity share in opposite commute directions.

\begin{figure}[H]
  \centering
  \safeincludegraphics[width=\textwidth]{illustrative_commute_overlay.pdf}
  \caption{Illustrative commute overlay applied to the base gravity share.
           The AM term increases residential-to-workplace OD weights, while
           the PM term reverses the direction. In the implementation these
           directional matrices are blended with the base share through
           Equation~\ref{eq:commute_blend} and then renormalised.}
  \label{fig:commute_overlay}
\end{figure}

\subsubsection{Stage 2 : Route library}
\label{sssec:route_library}

\paragraph{Edge expanded network}
\label{par:edge_expanded_network}

Routing for the scenario generator uses an \emph{edge-expanded} graph: SUMO
road edges are the graph nodes and a directed arc connects edge $u$ to edge
$v$ when SUMO allows the movement from $u$ to $v$. The weight of this arc is
the length of $v$, so a path cost corresponds to the sum of the edge lengths in
the route.

This representation is convenient for two reasons. First, the generator has to
sample SUMO edge endpoints and \texttt{duarouter} later expands those
endpoints into complete routes. Staying at edge level keeps the generated
demand close to SUMO's own trip format. Second, turning feasibility is already
encoded by SUMO connections: if an arc exists, the turn is allowed. Internal
junction edges\footnote{SUMO ids beginning with \texttt{:}} and edges that do
not permit private-vehicle access are excluded when the graph is built.


\paragraph{Impedance Estimation}
\label{par:impedance_estimation}

The gravity model needs a TAZ-to-TAZ impedance matrix. Computing exact
all-pairs shortest paths would be expensive and would add precision that is
not needed at this stage. Instead, for each directed zone pair $(i,j)$, the
generator samples a small number of origin-edge and destination-edge
combinations from the two zones. The lowest observed shortest-path cost is
used as the representative impedance. This is an approximation, but it gives
the gravity model the spatial signal it needs: nearby and well-connected zones
receive lower costs than distant or poorly connected ones.

\paragraph{$k$-Shortest Candidate Routes}
\label{par:k_shortest_candidate_routes}

For each active OD pair (i.e.\ $T^h_{ij} > 0$ for at least one hour $h$),
a set of at most $k$ candidate routes is computed as follows:
\begin{enumerate}
  \item Sample $s$ concrete (origin-edge, destination-edge)
        pairs, where edges are drawn proportionally to zone mass
        (Equation~\eqref{eq:zone_mass}).
  \item For each pair, enumerate up to $k$ shortest simple paths using
        Yen's algorithm \cite{yen_algorithm}
  \item Deduplicate across all samples by edge sequence, if the same
        sequence appears more than once, retain the lowest-cost
        instance.
  \item Return at most $k$ candidates sorted by cost.
\end{enumerate}

\subsubsection{Stage 3 : Trip Instantiation}
\label{sssec:vehicle_instantiation}

Once the hourly OD matrices and the route library are available, every integer
trip becomes a SUMO \texttt{trip} record. Two quantities are sampled: the
departure time and a feasible candidate path. Only the first and last edges of
that path are kept, because the final route is intentionally left to
\texttt{duarouter} in the Monte Carlo stage.

\paragraph{Departure Time}
\label{par:departure_time}

Trips within each one-hour interval $[b, b+3600)$ are assumed
to depart uniformly over that interval.  For an OD cell containing $n$ trips,
departure times are drawn independently as:
\begin{equation}
  t_1,\ldots,t_n \;\overset{\mathrm{i.i.d.}}{\sim}\; \mathcal{U}(b,\; b+3600).
\end{equation}
With only hourly totals available, this is the least committed assumption:
there is no information that would justify placing trips at a more precise
time within the hour.

\paragraph{Endpoint Selection}
\label{par:route_selection}

For an OD pair with candidate route set $\mathcal{R}_{ij}$, each trip first
selects one candidate path uniformly at random:
\begin{equation}
  P(r \mid \mathcal{R}) = \frac{1}{|\mathcal{R}|}, \quad r \in \mathcal{R}.
  \label{eq:uniform_choice}
\end{equation}
Only the first and last edges of the selected candidate are exported:
\begin{equation}
  \texttt{from} = r_1, \qquad \texttt{to} = r_{|r|}.
\end{equation}
This produces an unrouted \texttt{trips.xml} file. In this mode, the route
library is used only to choose feasible endpoints, it is not a fixed route
assignment. The full path choice is delegated to \texttt{duarouter} inside the
Monte Carlo loop, where edge costs are randomized independently for each run.

The departure-time and endpoint draws use one seeded random-number generator.
As a result, the generated \texttt{trips.xml} is reproducible. This matters
because the Monte Carlo experiment is meant to vary routing and simulation
seeds, not the underlying trip population.
